\section{Hypothesis: Stable ``Molecular Structure'' in Long CoT}

To understand these phenomena, as shown in Fig.~\ref{fig:main} \& \ref{fig:intro}, effective Long CoT connects nodes in logical space through reasoning behaviors, forming a stable macromolecular structure with mutually supportive components from a global perspective.

We formalize a Long CoT trace as a behavior-directed graph \(G = (V, E)\), where each node \(v \in V\) represents a reasoning step or edge \(s = (u, v) \in E\) is annotated with a behavior type \(s\rightarrow b \in \{\mathcal{D}, \mathcal{R}, \mathcal{E}\}\). For a trace corpus \(\mathcal{C}\), we estimate the behavior transition \(P_{\mathcal{C}}(b' \mid b)\) over consecutive edges and the marginal distribution \(\pi_{\mathcal{C}}(b)\).
Empirically, strong reasoning teachers produce stable \((P_{\mathcal{C}}, \pi_{\mathcal{C}})\) across models and tasks.
Specifically, these include the following three major bonds\footnote{Math definitions and proofs are in Appendix~\ref{append:math}.}:\vspace{-8pt}








\paragraph{\textbf{Deep Reasoning as Covalent Bonds}}
Deep reasoning forms the bone of the thought process, analogous to covalent bonds defining a molecule's primary chain. It encodes strong logical dependencies (Step A must justify Step B), maintaining direction and continuity; breaking this bone undermines the following steps and destabilizes the answer. By contrast, ``Normal Operation'' corresponds to stable local bonds within each step, capturing routine computation and direct semantic expression.\vspace{-8pt}

\paragraph{\textbf{Self-Reflection as Hydrogen Bonds}}
Reflection is a key stabilizer. As proteins gain stability when chains fold and form intra-chain hydrogen bonds, reasoning stabilizes when later steps (e.g., Step 100) test, revise, or reinforce earlier premises (e.g., Step 10). These long-range links constrain drift and hallucination, turning a long sequence into a more self-consistent structure. If later checks fail to align with earlier commitments, the reasoning cannot ``fold,'' indicating a structural logical error.


\paragraph{\textbf{Self-Exploration as Van der Waals Forces}}
Exploration resembles transient Van der Waals interactions: it supports abductive and inductive moves by enabling low-commitment associations in semantic space, where concepts can drift, combine, and be probed before stronger constraints are enforced.
